\documentclass{article}
\usepackage{PRIMEarxiv} % Core package for the PrimeAI Template
\bibliographystyle{plain}
\usepackage{amsmath, amssymb, amsthm, bm, dcolumn} % Add amsthm here for the proof environment
\usepackage[numbers,sort&compress]{natbib} % Natbib for citations
\usepackage{graphicx} % For high-quality images
\usepackage{multicol} % Multi-column figures/tables
\usepackage{hyperref} % Hyperlinks
\usepackage{listings} % Code listings
\usepackage{authblk} % For structured affiliations
% Define theorem style (optional)
\theoremstyle{plain}
\newtheorem{theorem}{Theorem}


% Custom commands for primes and qubit encoding
\newcommand{\primeQubit}[1]{|p_{#1}\rangle}
\newcommand{\primeGate}[1]{U_{p_{#1}}}

\usepackage{wrapfig}
\usepackage[pscoord]{eso-pic}
\usepackage[fulladjust]{marginnote}
\reversemarginpar

% Typesetting improvements without footnote patching
\usepackage[protrusion=true, expansion=true, tracking=false]{microtype}
\microtypecontext{spacing=nonfrench}

% Line numbers
\usepackage[right]{lineno}

% Text layout - adjust as needed
\raggedright
\setlength{\parindent}{0.5cm}
\textwidth 5.25in 
\textheight 8.75in

% Set double spacing
\usepackage{setspace} 
\doublespacing

% Adjust width for specific content
\usepackage{changepage}

% Adjust caption style
\usepackage[aboveskip=1pt,labelfont=bf,labelsep=period,singlelinecheck=off]{caption}

% Remove brackets from references
\makeatletter
\renewcommand{\@biblabel}[1]{\quad#1.}
\makeatother

% Header, footer, and page numbers
\usepackage{lastpage,fancyhdr}
\pagestyle{fancy}
\fancyhf{}
\fancyhead[L]{Preprint - PrimeAI Enhanced Template}
\fancyfoot[C]{\scriptsize Multiplicity Theory © 2024 Ryan Van Gelder - Citizen Gardens\\ Licensed Under MIT and CC BY-NC-SA 4.0.}
\fancyfoot[R]{Page \thepage\ of \pageref{LastPage}}
\renewcommand{\footrule}{\hrule height 2pt \vspace{2mm}}

\begin{document}
% Title and Author Info
\title{Integrating Grover's Search\\ With Multiplicity Theory}
\author{Ryan Van Gelder}
\affil{Citizen Gardens - The Foundation of Multiplicity \\ \texttt{ryan@citizengardens.org}}
\date{\today}

\maketitle

\begin{abstract}
Grover's search algorithm is a cornerstone of quantum computing, offering quadratic speedup for unstructured search problems~\cite{grover1996}. This paper expands Grover’s algorithm using Multiplicity Theory, incorporating prime encoding, recursive feedback, and quantum entanglement. These additions enhance search efficiency, cryptographic resilience, and adaptability. Furthermore, we explore the integration of Multiplicity’s prime-based structures into Grover’s framework, demonstrating applications in number theory, structured optimization, and quantum cryptography.
\end{abstract}

\section{Introduction}
Grover’s algorithm provides a quantum mechanical solution for unstructured search with a complexity of \(O(\sqrt{N})\) \cite{grover1996}. Multiplicity Theory, which uses prime encoding and recursive dynamics, complements Grover’s approach by improving state stability, adaptability, and scalability \cite{berkowitz2021}. This paper integrates Grover's algorithm with the principles of Multiplicity to optimize quantum search processes and expand its application in number theory and optimization problems.

\subsection{Motivation and Objectives}
This work aims to:
\begin{itemize}
    \item Extend Grover’s algorithm with prime-based encoding for compact and stable quantum states \citep{hardy1916}.
    \item Introduce recursive feedback to dynamically optimize Grover iterations in real-time \citep{childs2010}.
    \item Leverage entanglement among prime-encoded states for parallel computation \citep{ricks2004}.
\end{itemize}

\section{Methodology}

\subsection{Prime Encoding of Quantum States}
Prime encoding maps quantum states \(|\psi_i\rangle\) to primes \(\{p_1, p_2, \ldots, p_N\}\):
\begin{equation}
f(\psi_i) = p_i,
\end{equation}
where \(p_i\) is the prime corresponding to \(|\psi_i\rangle\) \citep{niven1991}. The superposition of prime-encoded states is given by:
\begin{equation}
|\Psi\rangle = \frac{1}{\sqrt{N}} \sum_{i=1}^{N} |p_i\rangle.
\end{equation}

\subsection{Prime-Encoded Oracle and Diffusion Operators}
The oracle identifies a target state \(p_{\text{target}}\) by flipping its phase:
\begin{equation}
O|p_i\rangle =
\begin{cases}
    -|p_i\rangle, & \text{if } p_i = p_{\text{target}},\\
    |p_i\rangle, & \text{otherwise.}
\end{cases}
\end{equation}
The diffusion operator amplifies the amplitude of \(|p_{\text{target}}\rangle\):
\begin{equation}
D = 2|\Psi\rangle\langle\Psi| - I.
\end{equation}

\subsection{Recursive Feedback Mechanism}
Dynamic adjustments optimize Grover iterations using a feedback mechanism:
\begin{equation}
k = g(k, t) = k_{\text{optimal}} + \alpha(t),
\end{equation}
where \(k_{\text{optimal}}\) is the optimal number of iterations and \(\alpha(t)\) adjusts for real-time conditions \citep{childs2010}.

\subsection{Entanglement of Prime-Encoded States}
Entanglement facilitates parallel computation. For two prime-encoded states \(|p_i\rangle\) and \(|p_j\rangle\), the entangled state is:
\begin{equation}
|\psi_{ij}\rangle = \frac{1}{\sqrt{2}} \left(|p_i\rangle|p_j\rangle + |p_j\rangle|p_i\rangle\right).
\end{equation}

\section{Results and Discussion}

\subsection{Applications in Number Theory}
Prime-encoded Grover’s search improves efficiency in:
\begin{itemize}
    \item \textbf{Prime Factorization:} Locating prime factors of composite numbers~\cite{shor1994}.
    \item \textbf{Modular Arithmetic:} Solving congruences of the form \(ax \equiv b \pmod{n}\).
    \item \textbf{Structured Optimization:} Identifying solutions constrained by prime properties~\cite{berkowitz2021}.
\end{itemize}

\subsection{Enhanced Probability of Success}
The success probability after \(k\) iterations is:
\begin{equation}
P(|p_{\text{target}}\rangle) \approx \sin^2\left((2k+1)\theta\right),
\end{equation}
where \(\theta\) is the amplitude rotation angle. Recursive feedback ensures optimal iteration counts~\cite{ricks2004}.

\subsection{Quantum Efficiency Gains}
Prime encoding reduces the effective search space to approximately \(\frac{N}{\ln N}\), enhancing computational efficiency~\cite{hardy1916}.

\section{Conclusion}
The integration of Lov Grover's algorithm with Multiplicity Theory creates a powerful synergy. By embedding prime encoding, recursive feedback, and entanglement, we achieve enhanced efficiency and adaptability in quantum search processes. Future work will explore multi-prime superpositions and hybrid frameworks integrating Shor’s and Grover’s algorithms~\cite{shor1994}.

\section*{Acknowledgments}
Special thanks to the Citizen Gardens Research Collective for their invaluable support.

\bibliographystyle{plain}
\bibliography{reference}

\end{document}
